% Part: sets-functions-relations
% Chapter: sets
% Section: unions-and-intersections

\documentclass[../../../include/open-logic-section]{subfiles}

\begin{document}

\olfileid[tr]{sfr}{set}{uni}
\olsection{Birleşimler ve Kesişimler}

\begin{explain}
\olref[sfr][set][bas]{sec} kesiminde soyutlama yoluyla, yani
$\Setabs{x}{\phi(x)}$ biçiminde küme tanımları vermiştik. Burada bir
özellik~$\phi$ kullanırız ve bu özellik, daha önce tanımladığımız kümelerden
söz edebilir. Örneğin $A$ ve~$B$ kümelerse
$\Setabs{x}{x \in A \lor x \in B}$ kümesi, $A$ kümesinin veya~$B$ kümesinin elemanı
olan bütün nesnelerden oluşur; başka bir deyişle $A$ ve~$B$ kümelerinin
elemanlarını bir araya getirir. Bunu, taralı bölgenin iki kümenin birlikte
alınan elemanlarını gösterdiği \olref{fig:union} ile görselleştirebiliriz.

\begin{figure}
  \olasset{assets/diagrams/union.tikz}
  \caption{İki kümenin birleşimi $A \cup B$, $A$ ve $B$ kümelerinin tüm
    elemanlarından oluşan kümedir.}
  \ollabel{fig:union}
\end{figure}

Kümeler üzerinde yapılan bu işlem---onları bir araya getirmek---çok yararlı ve
yaygındır; bu nedenle ona biçimsel bir ad ve simge veririz.
\end{explain}

\begin{defn}[Birleşim]
$A$ ve $B$ kümelerinin \emph{birleşimi}, $A \cup B$ ile gösterilen ve $A$
kümesinin, $B$ kümesinin ya da her ikisinin elemanı olan bütün nesnelerden
oluşan kümedir.
\[
A \cup B = \Setabs{x}{x \in A \lor x \in B}
\]
\end{defn}

\begin{ex}
Bir elemanın kaç kez yinelendiği önemli olmadığından, ortak bir elemanı
bulunan iki kümenin birleşiminde bu eleman yalnızca bir kez yer alır; örneğin
$\{ a, b, c\} \cup \{ a, 0, 1\} = \{a, b, c, 0, 1\}$.

Bir küme ile onun bir alt kümesinin birleşimi yalnızca büyük olan kümedir:
$\{a, b, c \} \cup \{a \} = \{a, b, c\}$.

Bir kümenin boş kümeyle birleşimi yine o kümedir:
$\{a, b, c \} \cup \emptyset = \{a, b, c \}$.
\end{ex}

\begin{prob}
$A \subseteq B$ ise $A \cup B = B$ olduğunu ispatlayın.
\end{prob}

\begin{explain}
Birleşimin \emph{duali} olan işlemi de ele alabiliriz. Bu işlem, hem $A$
kümesinin hem de $B$ kümesinin elemanı olan bütün nesnelerden oluşan kümeyi kurar. İşleme
\emph{kesişim} denir ve \olref{fig:intersection} ile gösterilebilir.
\begin{figure}
  \olasset{assets/diagrams/intersection.tikz}
  \caption{İki kümenin kesişimi $A \cap B$, $A$ ve $B$ kümelerinin ortak
    elemanlarından oluşan kümedir.}
  \ollabel{fig:intersection}
\end{figure}
\end{explain}

\begin{defn}[Kesişim]
$A$ ve $B$ kümelerinin \emph{kesişimi}, $A \cap B$ ile gösterilen ve hem
$A$ kümesinin hem de~$B$ kümesinin elemanı olan bütün nesnelerden oluşan
kümedir.
\[
A \cap B = \Setabs{x}{x \in A \land x \in B}
\]
Kesişimleri boş olan iki kümeye \emph{ayrık} denir. Bu, ortak elemanlarının
bulunmadığı anlamına gelir.
\end{defn}

\begin{ex}
İki kümenin hiçbir ortak elemanı yoksa kesişimleri boştur:
$\{ a, b, c\} \cap \{ 0, 1\} = \emptyset$.

İki kümenin ortak elemanları varsa kesişimleri bu elemanların tümünden oluşan
kümedir:
$\{a, b, c \} \cap \{a, b, d \} = \{a, b\}$.

Bir küme ile onun bir alt kümesinin kesişimi yalnızca küçük olan kümedir:
$\{a, b, c\} \cap \{a, b\} = \{a, b\}$.

Herhangi bir kümenin boş kümeyle kesişimi boştur:
$\{a, b, c \} \cap \emptyset = \emptyset$.
\end{ex}

\begin{prob}
$A \subseteq B$ ise $A \cap B = A$ olduğunu titizlikle ispatlayın.
\end{prob}

\begin{explain}
İkiden fazla kümenin birleşimini veya kesişimini de oluşturabiliriz. Bunu genel
olarak ele almanın zarif bir yolu şudur: birleşimini (veya kesişimini)
oluşturmak istediğimiz bütün kümeleri tek bir kümede topladığımızı varsayalım.
Bu durumda başlangıçtaki kümelerin birleşimini, topladığımız kümenin elemanı
olan kümelerden en az birinin elemanı olan bütün nesnelerin kümesi; kesişimini
ise bu kümelerin her birinin elemanı olan bütün nesnelerin kümesi olarak
tanımlayabiliriz.
\end{explain}

\begin{defn}
$A$ bir kümeler kümesiyse $\bigcup A$, $A$ kümesinin elemanları olan
kümelerin elemanlarından oluşan kümedir:
\begin{align*}
\bigcup A & = \Setabs{x}{x \text{, } A\text{ kümesindeki bir kümeye aittir}},
\text{ yani}\\
& = \Setabs{x}{\text{öyle bir } B \in A
  \text{ vardır ki } x \in B}
\end{align*}
\end{defn}

\begin{defn}
$A$ bir kümeler kümesiyse $\bigcap A$, $A$ kümesinin elemanı olan her kümede
bulunan nesnelerden oluşan kümedir:
\begin{align*}
\bigcap A & = \Setabs{x}{x \text{, } A\text{ kümesindeki her kümeye aittir}},
\text{ yani}\\
 & = \Setabs{x}{\text{her } B \in A \text{ için } x \in B}
\end{align*}
\end{defn}

\begin{ex}
$A = \{ \{ a, b \}, \{ a, d, e \}, \{ a, d \} \}$ olduğunu varsayalım.
Bu durumda $\bigcup A = \{ a, b, d, e \}$ ve $\bigcap A = \{ a \}$ olur.
\end{ex}
\begin{prob}
	$A$ bir kümeyse ve $A \in B$ ise $A \subseteq \bigcup B$ olduğunu gösterin.
\end{prob}

Aynı işlemleri $A_1$, $A_2$, \dots{} küme dizisi için de yapabiliriz:
\begin{align*}
\bigcup_i A_i & = \Setabs{x}{x \text{, } A_i\text{ kümelerinden birine aittir}}\\
\bigcap_i A_i & = \Setabs{x}{x \text{, her } A_i\text{ kümesine aittir}}.
\end{align*}

Kümeleri indislemek üzere, her $i \in I$ için $A_i$ kümesini ele aldığımız bir
$I$ kümesi verilmişse şu kısaltmaları da kullanabiliriz:
\begin{align*}
	\bigcup_{i \in I} A_i & = \bigcup \Setabs{A_i }{i \in I}\\
	\bigcap_{i \in I} A_i & = \bigcap\Setabs{A_i}{i \in I}
\end{align*}

Son olarak $A$ kümesinde bulunup $B$ kümesinde bulunmayan bütün elemanlardan
oluşan kümeyi düşünmek isteyebiliriz. Bunu \olref{difference} ile
gösterebiliriz.

\begin{figure}
  \olasset{assets/diagrams/difference.tikz}
  \caption{$A$ kümesinin $B$ kümesinden farkı $A \setminus B$, $A$
    kümesinde bulunup $B$ kümesinde bulunmayan elemanlardan oluşur.}
  \ollabel{difference}
\end{figure}

\begin{defn}[Küme Farkı]
$A$ kümesinin $B$ kümesinden \emph{farkı}~$A \setminus B$, $A$'da bulunup
$B$'de bulunmayan bütün elemanlardan oluşur; yani
\[
A\setminus B = \Setabs{x}{x\in A \text{ ve } x \notin B}.
\]
\end{defn}

\begin{prob}
	$A \subsetneq B$ ise $B \setminus A \neq \emptyset$ olduğunu ispatlayın.
\end{prob}

\end{document}
